% Generated draft: requires mathematical and source review.
\documentclass[11pt]{article}
\usepackage[margin=25mm]{geometry}
\usepackage{fontspec}
\setmainfont{Latin Modern Roman}
\usepackage{amsmath,amssymb,mathtools,unicode-math}
\setmathfont{Latin Modern Math}
\usepackage{enumitem,needspace,xurl,hyperref}
\setlength{\parindent}{0pt}
\setlength{\parskip}{5pt}
\setlength{\emergencystretch}{3em}
\newcommand{\N}{\mathbb{N}}
\newcommand{\Z}{\mathbb{Z}}
\newcommand{\Q}{\mathbb{Q}}
\newcommand{\R}{\mathbb{R}}
\newcommand{\C}{\mathbb{C}}
\newcommand{\F}{\mathbb{F}}
\newcommand{\A}{\mathbb{A}}
\newcommand{\PP}{\mathbb{P}}
\newcommand{\E}{\mathbb{E}}
\newcommand{\eps}{\varepsilon}
\newcommand{\dd}{\,\mathrm{d}}
\begin{document}
\section*{UnsolvedMath ID: 800505}
\textbf{Source problem code:} catalogue entry
\subsection{Definitions and mathematical statement}
A conventional framework for expressing the question is the following. Let \(\mathbb K\in\{\R,\C\}\), let \(A\in\mathbb K^{n\times n}\), and consider \(Ax=b\) with initial approximation \(x_0\) and initial residual \(r_0=b-Ax_0\). Define
\[
\mathcal K_0(A,r_0)=\{0\},\qquad
\mathcal K_k(A,r_0)=\operatorname{span}\{r_0,Ar_0,\ldots,A^{k-1}r_0\}\quad(k\geq1).
\]
For standard unrestarted GMRES in the Euclidean norm, the residual \(r_k=b-Ax_k\) satisfies
\[
x_k\in x_0+\mathcal K_k(A,r_0),\qquad
\|r_k\|_2=\min_{x\in x_0+\mathcal K_k(A,r_0)}\|b-Ax\|_2.
\]
These conventions provide notation; the supplied record does not specify the GMRES variant, residual norm, or precise discretized system in the examples it refers to.

Eigenvalue information means the locations of the elements of \(\sigma(A)\), the spectrum of \(A\). If \(A\) admits a diagonalization \(A=V\Lambda V^{-1}\), eigenvector conditioning can be expressed, for a specified eigenvector matrix \(V\), by
\[
\kappa_2(V)=\|V\|_2\,\|V^{-1}\|_2.
\]
This conditional definition does not assume that the convection--diffusion matrices in the source are diagonalizable. A choice of eigenvector normalization and, for repeated eigenvalues, a choice of eigenbasis would be needed to make such a measure unambiguous.

\textbf{Problem (supplied source formulation).} Explain the acceleration of GMRES convergence after an initial transient phase observed in some convection--diffusion model problems, including the example cited in the source as [LieStr05]. In particular, determine whether, and in what precise sense, the convergence speed can be related to the eigenvalue locations and eigenvector conditioning of the discretized convection--diffusion operators.

The supplied formulation is an explanatory research question, rather than a quantified assertion of a particular convergence estimate. It does not define the transient phase or acceleration, identify the differential operators, domains, coefficients, boundary conditions, discretizations, right-hand sides, initial guesses, or preconditioners, or specify which spectral data should predict which measure of convergence. The bibliographic identity and model problem of [LieStr05] are also absent. Consequently, no universal quantifier over convection--diffusion systems, particular onset criterion, or specific convergence bound is asserted here; those unresolved qualifications require examination of the cited source.
\subsection{Short English statement}
Explain why GMRES sometimes converges faster after an initial transient when applied to discretized convection--diffusion problems, especially the example referenced as [LieStr05]. Determine whether eigenvalue locations and eigenvector conditioning explain that behavior. The supplied record leaves the model, algorithmic details, and meaning of acceleration unspecified.
\subsection{Sources}
\begin{itemize}
\item Local UnsolvedMath registry, permanent ID \(800505\), source problem code \texttt{catalogue entry}, title ``Linear Systems and Eigenvalue Problems --- Problem 14.'' The supplied record attributes the problem to Amsel, Baumann et al. (2026) and labels its status as open. This is a supplied registry label, not an independently verified mathematical or literature assessment. The supplied upstream revision is \texttt{d371ee917ac9647fbb06cd6fc529deec6f3dda82}.
\item Amsel, Baumann et al., \emph{Linear Systems and Eigenvalue Problems: Open Questions from a Simons Workshop} (2026), Problem 14. Supplied source URL: \url{https://arxiv.org/abs/2602.05394}. The record reports access on 2026-10-16 and extraction by \texttt{arxiv\_latex}; these are supplied provenance details, not a claim of independent access or verification.
\item [LieStr05], cited in the supplied statement for a convection--diffusion model problem exhibiting the reported behavior. No authors, title, publication details, or URL for this reference are supplied, so its identity and the precise example remain unresolved.
\end{itemize}
\end{document}
